% tbsrc/intro.tex — Introduction (Section 1).  STATUS: written (first draft).
\section{Introduction}
\label{sec:intro}

\subsection{Positron emission tomography in one page}
\label{sec:intro:pet}

Positron Emission Tomography (PET) images the spatial distribution of a
radiotracer labelled with a positron-emitting isotope. An emitted positron
travels a short distance (the \emph{positron range}) before annihilating with
an atomic electron; the annihilation converts the pair's rest mass into two
$511\,\mathrm{keV}$ photons emitted almost exactly back to back. A ring of
scintillation detectors surrounds the subject and records pairs of photons that
arrive within a narrow coincidence time window. Each such \emph{coincidence}
is assigned to the straight line joining the two firing detector elements ---
the \emph{line of response} (LOR). The annihilation is known to have occurred
somewhere along that line (in non-time-of-flight PET, with no further
localization along it).

The raw output of the scanner is therefore a large set of LORs, each carrying
the number of coincidences recorded on it. The data can be stored as a
\emph{sinogram} (counts binned by LOR), or, event by event, as a
\emph{listmode} stream in which every coincidence is one record giving the two
endpoint positions. This tutorial works primarily in listmode, which is the
natural representation for the continuous, machine-learning--positioned
detectors that motivate \code{RecoCrysp} (Section~\ref{sec:resolution}).

\subsection{The reconstruction problem}
\label{sec:intro:problem}

Let $f(\bm{u})$ be the activity concentration at world position
$\bm{u}\in\RR^3$. The expected number of (true, unscattered) coincidences on a
LOR is, to first approximation, the line integral of $f$ along it. Discretizing
$f$ on a voxel grid (Section~\ref{sec:grid}) turns it into a vector
$\xv\in\RR^{N}_{\ge 0}$, and the line-integral operation becomes a linear map
--- the \emph{system matrix} or \emph{forward projector} $\Aop$ --- so that the
mean data over a set of LORs is, in the idealized case,
\begin{equation}
  \pred \;=\; \Aop\,\xv .
  \label{eq:intro-ideal}
\end{equation}
Reconstruction is the inverse problem: recover $\xv$ from the measured counts
$\yv$. It is \emph{ill-posed} --- the data are noisy (finite counts), angularly
and radially incompletely sampled, and blurred by the finite detector
resolution and physical effects --- so a naive inversion of \eqref{eq:intro-ideal}
is useless. Instead one adopts a statistical model of the data and seeks the
image that best explains them. Because PET counts are Poisson distributed
(Section~\ref{sec:statistics}), the method of choice is
\emph{maximum-likelihood expectation-maximization} (MLEM) and its accelerated
ordered-subsets variant (OSEM), developed in
Sections~\ref{sec:mlem}--\ref{sec:acceleration}.

The realistic forward model is richer than \eqref{eq:intro-ideal}. Writing the
expected counts on LOR $i$,
\begin{equation}
  \pred_i \;=\; n_i\, a_i \,\bigl(\Aop\,\Gop\,\xv\bigr)_i \;+\; s_i + r_i ,
  \label{eq:intro-full}
\end{equation}
collects every effect this tutorial treats: $\Gop$ is the detector
point-spread function (Section~\ref{sec:resolution}); $n_i$ the detector
\emph{normalization} and $a_i$ the \emph{attenuation} of the medium (both
multiplicative, Section~\ref{sec:datamodel}); and $s_i,r_i$ the additive
\emph{scatter} and \emph{randoms} backgrounds. Each tutorial example switches
on a subset of these terms (Section~\ref{sec:summary}); the first example sets
$\Gop$ to the detector resolution, $a=1$ (emission in vacuum), $n=1$, and
$s=r=0$.

\subsection{Scope, library, and notation}
\label{sec:intro:scope}

This document is the shared \emph{basis} for the tutorial series: it states the
physics and mathematics once, so the example documents can focus on a concrete
setup and its results. The treatment is non-time-of-flight throughout. All
computations are carried out with \code{RecoCrysp}, a Julia library whose
projectors are written once with \code{KernelAbstractions.jl} and run
unchanged on multithreaded CPUs and on GPUs, including Apple Metal; all
arithmetic is single precision (\code{Float32}). The library is endpoint-driven
--- it knows only LORs, voxels, and the operators of
\eqref{eq:intro-full} --- which is what makes it agnostic to the detector
technology (Section~\ref{sec:geometry}).

Notation used throughout is collected in Appendix~\ref{app:notation}; the main
symbols are summarized in Table~\ref{tab:intro-notation}.

\begin{table}[t]
  \centering
  \begin{tabular}{ll}
    \toprule
    Symbol & Meaning \\
    \midrule
    $\xv\in\RR^N_{\ge0}$ & activity image (voxel vector) \\
    $\yv$               & measured counts (per LOR / per event) \\
    $\pred$             & expected (mean) counts, Eq.~\eqref{eq:intro-full} \\
    $\Aop,\ \Aadj$      & forward projector and its matched adjoint \\
    $\Gop$              & detector point-spread (resolution) operator \\
    $n_i,\ a_i$         & normalization, attenuation (multiplicative) \\
    $s_i,\ r_i$         & scatter, randoms (additive) \\
    $\sensv=\Aadj(\nrm\!\cdot\!\att)$ & sensitivity image \\
    \bottomrule
  \end{tabular}
  \caption{Principal symbols. Per-LOR factors are written without indices when
           viewed as vectors.}
  \label{tab:intro-notation}
\end{table}
